\documentclass[10pt,a4paper]{amsart}
\usepackage[margin=19mm]{geometry}
\usepackage{amsmath,amssymb,mathtools}
\usepackage[scr=rsfs,cal=euler]{mathalfa}
\usepackage{xfrac}
\usepackage[unicode,hidelinks,bookmarks=false]{hyperref}
\newcommand{\ie}{i.e.\ }
\newcommand{\cf}{cf.\ }
\newcommand{\resp}{respectively\ }
\newcommand{\Subsection}{\subsection}
\providecommand{\nobreakdash}{\nobreak\hbox{-}}
\setlength{\emergencystretch}{2em}
\multlinegap0pt
\usepackage{amssymb}
\usepackage[shortlabels]{enumitem}
\usepackage{float}
\usepackage{mathtools}
\usepackage{tikz}
\newtheorem{theorem}{Theorem}
\newtheorem{lemma}{Lemma}
\newcommand{\B}[1]{{\mathbf #1}}
\title{On quasimorphisms and distortion in homeomorphism groups}
\author{Michael Brandenbursky, Jarek Kedra, Michal Marcinkowski, Egor Shelukhin}
\date{Source: arXiv:2512.15375v2 (CC BY 4.0)}
\begin{document}
\maketitle

\noindent\emph{Source and licence.} On quasimorphisms and distortion in homeomorphism groups. Version arXiv:2512.15375v2 (CC BY 4.0), distributed by arXiv under the Creative Commons Attribution 4.0 International licence, http://creativecommons.org/licenses/by/4.0/, as recorded in the arXiv licence metadata for this exact version and shown on the abstract page https://arxiv.org/abs/2512.15375v2. Full licence notice: \texttt{licenses/arxiv-2512.15375-CC-BY-4.0.txt}.\par\smallskip

\noindent\textit{Editorial scope.} Counted unit: the article's theorem T:distortion (source0.tex lines 339--346). The article's complete original proof of it is delivered with it (source0.tex lines 348--366, one \texttt{proof} environment of 19 lines). No mathematics is written, repaired or re-derived: the statement and the proof are the article's own lines, in the article's own notation, and no retained line is substituted or re-flowed.  Retained uncounted as the source-local context the proof needs: lines 529--537.  Nothing was omitted from any retained span, so the package carries no omission record.  No retained line contains a citation, so the wrapper carries no bibliography and no bibliography entry.  No retained line contains a figure environment, an image-inclusion command or drawing code, so no image is delivered and the package declares no asset dependency.  Editorial formatting only, in addition: an AMS \texttt{amsart} preamble that restates the article's own declarations and macros used by the retained lines, namely the environments \texttt{theorem}, \texttt{lemma} on separate counters, together with the standard packages those lines need.  The retained environments print in this document as Theorem 1, Lemma 1 in order of appearance, whereas the article prints the same items with the numbering of its own full document.  The complete native source of the article as distributed in the arXiv e-print for this exact version is bundled unchanged as \texttt{sources/191167/source0.tex}.
\begin{lemma}\label{L:psi-norm}
Let $\delta\psi$ be a $1$-bounded cocycle.
The formula
$$
|f|_{\psi} = \sup_g|\psi(g)-\psi(fg)|
$$
defines a pseudo-norm on $G$. Moreover, $\psi$ is $1$-Lipschitz with
respect to the pseudo-metric defined by $d_{\psi}(f,g) = |fg^{-1}|_{\psi}$.
\end{lemma}

\begin{theorem}\label{T:distortion}
Let $\psi \colon G\to \B R$ be a function such that $\delta\psi$ is semi-bounded.
If 
$$
\limsup_{k\to \infty}\frac{|\psi(f^k)|}{k}>0
$$ 
then $f$ is undistorted.
\end{theorem}

\begin{proof}
It follows from Lemma \ref{L:psi-norm} that $|f|_{\psi}=\sup_g|\psi(g)-\psi(fg)|$ 
is a pseudo-norm with respect to which
$\psi$ is Lipschitz which implies that 
\begin{equation}
\limsup_{k\to \infty} \frac{|f^k|_{\psi}}{k}>0.
\label{Eq:limsup}
\end{equation}
Let
$\Gamma\subseteq G$ be a finitely generated subgroup containing $f$
and let $C=\max\{|s|_{\psi}\ |\ s\in S\}$, 
where $S\subseteq \Gamma$ is a finite generating set. Then
$$
|f|_{\psi} = |s_1\dots s_m|_{\psi} \leq |s_1|_{\psi} + \dots + |s_m|_{\psi}
\leq C |f|_S
$$
which, together with \eqref{Eq:limsup} implies that $f$ is undistorted in $\Gamma$
and hence in $G$.
\end{proof}

\end{document}
